% Statement and cited-input packet. Target proof is separately embargoed.
% Frozen II source SHA256: 74497bfe327fae25c9cdcb17334b3bf42a9793a9df7c0a4404d97568f2f760f0
% Target owner: push:src:E.2.6; tag: 033E

\begin{proposition}[positive canonical slope preserves the fibration]\hypertarget{tag-033E}{}\label{push:src:E.2.6}


Under the hypotheses of \ref{push:src:E.2.5}, suppose $D>0$.  Then $E_0=0$
and $e=1$.  After shrinking the disc and choosing a frame of
$\pi_*\mathcal L$, its evaluation defines a proper flat
holomorphic map
\[
 g:\mathcal X\longrightarrow\mathbb P^1\times\Delta
\]
which is the intrinsic algebraic reduction on every fibre.
Its smooth general central fibre is a complex two-torus.

In particular, two families of the type in \ref{push:src:E.2.5} with positive
canonical slope cannot have the same smooth central manifold
unless their multiplicity pairs agree.  A positive-slope and
a negative-slope family cannot have a common smooth central
manifold either.  This proposition concerns a common smooth
specialization.  The additional propagation and closure
arguments for arbitrary connected chains are \ref{push:src:E.2.22}--\ref{push:src:E.2.24}.
\end{proposition}

% Cited input push:src:E.2.22; statement or labelled formula only.
\begin{proposition}[propagation from a known open subset of a disc]\hypertarget{tag-0338}{}\label{push:src:E.2.22}


Let $\pi:\mathcal X\to\Delta$ be a proper holomorphic
submersion with fibres diffeomorphic to $S^6$.  Suppose
there is a nonempty open subset $U\subset\Delta$ such
that every $X_s$, $s\in U$, is a member of the fixed
coprime multiplicity family $(3k,4l)$ of \SeriesRef{push:src:E.2.1}.
Then every fibre over $\Delta$ is a member of that
same family.
\end{proposition}

% Cited input push:src:E.2.24; statement or labelled formula only.
\begin{corollary}[separation for continuous complex structures]\hypertarget{tag-033A}{}\label{push:src:E.2.24}


Let $\mathcal I(S^6)$ be the space of smooth integrable
almost-complex structures on the fixed smooth sphere,
with its $C^\infty$ topology.  For a fixed coprime pair
$(3k,4l)$, let $\mathcal I_{k,l}$ consist of the structures
whose complex manifolds are biholomorphic to constructed
members with that pair.  Then $\mathcal I_{k,l}$ is open
and closed in $\mathcal I(S^6)$.

Consequently a continuous family of such structures
over any connected topological parameter space cannot
contain constructed members with two different pairs.
The same separation holds in the quotient spaces by
$\operatorname{Diff}_0(S^6)$ or by all orientation-preserving
diffeomorphisms, with their quotient topologies.
This statement does not identify the number of marked
connected components belonging to a single pair.
\end{corollary}

% Cited input push:src:E.2.5; statement or labelled formula only.
\begin{proposition}[a fractional canonical pencil at a smooth limit]\hypertarget{tag-033D}{}\label{push:src:E.2.5}


Let $\pi:\mathcal X\to\Delta$ be a proper holomorphic submersion
whose fibres are diffeomorphic to $S^6$.  Suppose that for every
$t\ne0$ the fibre is one of the members of \SeriesRef{push:src:E.2.1}, with the same
coprime multiplicities $(3k,4l)$.  Set
\[
 A=12kl,\qquad D=kl-3k-4l\ne0.
\]
There is a unique holomorphic line bundle $\mathcal L$ on
$\mathcal X$, up to isomorphism, satisfying
\[
 \mathcal L^{D}=K_{\mathcal X/\Delta}^{A}.
\]
Its restriction to $X_t$ for $t\ne0$ is the coarse pencil
$f_t^*\mathcal O_{\mathbb P^1}(1)$.  The coherent sheaf
$\pi_*\mathcal L$ is locally free of rank two, and its central
fibre injects into $H^0(X_0,\mathcal L_0)$.

The central manifold has algebraic dimension one and a
holomorphic algebraic reduction $g_0:X_0\to\mathbb P^1$ with
connected fibres.  If $s_0,s_1$ are the two central sections
obtained from a local frame of $\pi_*\mathcal L$, there are a
nonconstant rational map $R:\mathbb P^1\to\mathbb P^1$ of
degree $e\ge1$ and an effective Cartier divisor $E_0$ such that
\[
 [s_0:s_1]=R\circ g_0,\qquad
 \mathcal L_0\simeq g_0^*\mathcal O(e)\otimes\mathcal O(E_0).
\]
Here $[s_0:s_1]$ denotes the holomorphic extension of the ratio;
the displayed divisor is its common fixed divisor.  The
statement does not assert that $E_0=0$.
\end{proposition}

